\section{Lean correspondence}\label{app:lean}

The table lists, for each statement of the paper with machine-checked content, the Lean theorems and their files, relative to the directory \texttt{lean/} that accompanies the paper. A row marked ``by hand'' has no Lean counterpart. The notation of the Lean sources and of their comments differs from the paper's as follows (\texttt{lean/README.md} has the same list as a table). The field is \texttt{K}, with \texttt{q} elements, for $\E$ (except in \texttt{Interleave}, where \texttt{K} is the small field and \texttt{F} or \texttt{E} an extension, and in \lean{protocolD_sound_of_HK}, where the small field is \texttt{F}). A word is \texttt{O}, with halves \texttt{O}\textsubscript{\texttt{0}}, \texttt{O}\textsubscript{\texttt{1}} and fold \texttt{O}\textsubscript{\texttt{c}}, for $\orc$, $\orc_0$, $\orc_1$, $\orc_c$. The degree bound is \texttt{d} for $\dgr$. An agreement is \texttt{A} for $\agr$, and the thresholds are \texttt{thr k}, written $A^{(i,k)}$ in the comments, for $\agr^{(i,k)}$. A fold challenge is \texttt{c = (c.1, c.2)}, written $(\alpha,\beta)$ in the comments, for $(\afa,\afb)$. The running claim's value is $\sigma$ for $c$. The prover's scalars are \texttt{vG}, \texttt{vF}, $\sigma_0$, $\sigma_1$, $\varphi$\texttt{m} for $v_G,v_F,\tabx{0},\tabx{1},\varphi_m$. The number of queries is \texttt{t} for $\nq$. The discount is $\tau$ in both. The threshold of the triple count is \texttt{B} for $\thr=\Kz(\Kz-1)(\Kz-M)$. The bound on the lists of the folds is $\Lambda$ where the paper writes $L$; it appears in the verbatim listing of \lean{lemmaF_exc} below. And \texttt{dens p}, the fraction of the elements of a finite type that satisfy \texttt{p}, is a probability over a uniform verifier message.

{\small
\begin{longtable}{@{}L{0.25\textwidth}L{0.42\textwidth}L{0.27\textwidth}@{}}
\toprule
statement & Lean theorems & file under \texttt{lean/}\\
\midrule
\endhead
Lemma~\ref{lem:paths} (paths), Lemma~\ref{lem:lines} & by hand & \\
\addlinespace
Proposition~\ref{prop:flats} (flats) & \lean{subPoly_isLinearized}, \lean{natDegree_subPoly_sub_le}, \lean{subPoly_adjoin} & \lean{Constructions/Subspace.lean}\\
 & \lean{flatPoly_natDegree_le}, \lean{flatPoly_eval}, \lean{agree_flatPoly}, \lean{flatPoly_injective}, \lean{flats_list}, \lean{card_hyperplanes_ge}, \lean{flats_quarter} & \lean{Constructions/Flats.lean}\\
 & \lean{flats_quarter_GF128} (an instance at $\blen=2^{26}$) & \lean{Constructions/Concrete.lean}\\
\addlinespace
Proposition~\ref{prop:five8}(a) (word $\loc_U^5$) & \lean{glued_list}, \lean{glued_five_eighths}, \lean{glued_eq_imp_invariant} & \lean{Constructions/Glued.lean}\\
 & \lean{glued_list_card} (at least $\blen/8-1$ members; the count $56\blen-392$ is by hand) & \lean{Constructions/GluedCount.lean}\\
Proposition~\ref{prop:five8}(b) & by hand & \\
Proposition~\ref{prop:five8}(c) (padded word) & \lean{padded_flats_list}, \lean{five_eighths_slack}, \lean{five_eighths_genuine} & \lean{Constructions/Padded.lean}\\
 & \lean{gap_lemma_locator}; \lean{flatUnion_five_eighths} (every union of flats of the family shares the top of its word) & \lean{Constructions/Gap.lean}, \lean{Constructions/UnionOfFlats.lean}\\
\addlinespace
Proposition~\ref{prop:pole} (pole collapse) & \lean{pole_collapse}, \lean{combiner_collapse}, \lean{affFold_pole_close_iff}, \lean{polyFold_pole_close_iff}, \lean{card_agree_poleWord_le}, \lean{lemmaA}, \lean{lemmaA_card}, \lean{list_gives_bad_challenges} & \lean{Constructions/Pole.lean}\\
 & \lean{flats_break_mca}, \lean{flats_break_mca_GF128} & \lean{Constructions/Concrete.lean}\\
\addlinespace
Proposition~\ref{prop:complete} (completeness) & by hand; uses the identities of Lemma~\ref{lem:E} & \\
\addlinespace
Lemma~\ref{lem:E} (evaluation functional) & \lean{evL_add}, \lean{evL_mul}, \lean{evL_halves}, \lean{card_evL_zero_le} & \lean{DeepFold/Lanes.lean}\\
 & \lean{exL_injective}, \lean{exL_add_le}, \lean{evL_exL_halves} & \lean{Protocol/Level.lean}\\
\addlinespace
Lemma~\ref{lem:L1} (determinant) & \lean{Delta_eq_det}, \lean{Delta_char_two}, \lean{Delta_eval_eq_zero}, \lean{Delta_natDegree_le} & \lean{DeepFold/Basic.lean}\\
 & \lean{isFold_unique}, \lean{Delta_eq_zero_iff} & \lean{DeepFold/Class.lean}\\
 & \lean{isFoldL_unique}, \lean{DeltaL_eq_zero_iff}, \lean{evL_DeltaL_eq_zero} & \lean{DeepFold/Lanes.lean}\\
 & \lean{card_nondep_triples_ge} & \lean{DeepFold/Counting.lean}\\
\addlinespace
Lemma~\ref{lem:L2} (class trichotomy) & \lean{class_trichotomy}, \lean{class_poly_of_card_gt} & \lean{DeepFold/Class.lean}\\
 & \lean{card_agr_outside_le}, \lean{agrPair_card_ge} & \lean{DeepFold/Agreement.lean}\\
 & \lean{classL_poly_of_card_gt}, \lean{card_agrL_outside_le}, \lean{exists_nearL} & \lean{DeepFold/Lanes.lean}\\
\addlinespace
Theorem~\ref{thm:count} (counting) & \lean{card_Exc_mul_le}, \lean{density_Exc_le}, \lean{exists_big_class} & \lean{DeepFold/Counting.lean}\\
 & \lean{card_ExcL_mul_le}, \lean{density_ExcL_le}, \lean{deepFoldL_theorem3} & \lean{DeepFold/Lanes.lean}\\
\addlinespace
Lemma~\ref{lem:F} (DEEP fold round) & \lean{lemmaF_exc}, \lean{lemmaF_round} & \lean{DeepFold/Main.lean}\\
 & \lean{card_quadric_zeros_le}, \lean{deepFold_lemmaF} & \lean{DeepFold/Claim.lean}\\
 & \lean{mem_fullT} (the set $\cand$) & \lean{DeepFold/Full.lean}\\
\addlinespace
Lemma~\ref{lem:P} (pinning) & \lean{lemmaP}, \lean{pinned_unique}, \lean{card_pinned_le_one} & \lean{Protocol/Pinning.lean}\\
\addlinespace
Lemma~\ref{lem:T} (potential) & \lean{pot_step}, \lean{potT_step}, \lean{potT_init} & \lean{Protocol/Potential.lean}\\
 & \lean{levelPot_init}, \lean{levelPot_step} (one batch) & \lean{Protocol/Level.lean}\\
 & \lean{ladderPot_init}, \lean{ladderPot_step} (a schedule) & \lean{Protocol/Ladder.lean}\\
\addlinespace
Theorem~\ref{thm:rbr-deep}, row M0 & \lean{listBound_level}, \lean{dens_not_good_le} & \lean{Protocol/Level.lean}\\
 & \lean{dens_not_goodS_le} & \lean{Protocol/Ladder.lean}\\
row M1 & \lean{batch_transition}, \lean{batch_transition_density} & \lean{Rounds/Batch.lean}\\
row M2 & \lean{fold_transition}, \lean{fold_transition_plain}, \lean{fold_transition_density} & \lean{Rounds/Fold.lean}\\
row M3 & \lean{query_transition}, \lean{query_transition_density} & \lean{Rounds/Query.lean}\\
row M4, row M5 & \lean{final_transition}, \lean{final_transition_plain}, \lean{final_transition_density}, \lean{tail_round_transition}, \lean{tail_round_transition_density} & \lean{Rounds/Final.lean}\\
M1, M3 from Lemma~\ref{lem:P} & \lean{batch_transition_pinned}, \lean{query_transition_pinned} & \lean{Rounds/Pinned.lean}\\
the definition of round-by-round soundness (\S\ref{sec:annexb}) and the union bound over a run & \lean{PublicCoin}, \lean{StateFunction}, \lean{RBRSound}, \lean{StateFunction.winProb_le}, \lean{RBRSound.winProb_le} & \lean{Soundness/RBR.lean}\\
the state of \S\ref{sec:state} with its potential, and each transition from that state & \lean{FullState}, \lean{points_full}, \lean{batch_full}, \lean{fold_full}, \lean{query_full}, \lean{final_full} & \lean{Soundness/State.lean}\\
Protocol~\ref{prot:D} with one batch per level, as an abstract model; Theorem~\ref{thm:rbr-deep} for it (properties (i) to (iv), rows M0 to M5) & \lean{protocolD}, \lean{Level.OK}, \lean{doomed}, \lean{err}, \lean{win}, \lean{step_tail}, \lean{protocolD_stateFunction}, \lean{protocolD_rbr} & \lean{Soundness/ProtocolD.lean}\\
Theorem~\ref{thm:rbr-deep} and Theorem~\ref{thm:K} over a run of the model: every strategy wins with probability at most the sum of the errors & \lean{budget}, \lean{protocolD_sound}, \lean{protocolD_sound_plain}, \lean{protocolD_sound_extract}, \lean{protocolD_sound_one_level} & \lean{Soundness/Sound.lean}\\
prover strategies exist on every ladder and with every tail, under two conditions on the abstract maps (\S\ref{sec:composed}) & \lean{Terminal.AdmitsZero}, \lean{Level.Fits}, \lean{ladderFits}, \lean{strategy_nonempty_of_fits}, \lean{strategy_nonempty_ladder}, \lean{strategy_nonempty_one_level}; \lean{finalMsg_isEmpty}, \lean{handover_isEmpty} (without the conditions the prover can have no message) & \lean{Soundness/Sound.lean}\\
the same from Hypothesis~\ref{hyp:HK}; with no list hypothesis in the unique-decoding range; two concrete instances with every hypothesis discharged (one level with no tail; two levels with a tail), neither a deployed parameter set; level-$0$ errors & \lean{Level.OK.of_rsListBound}, \lean{Level.HK}, \lean{protocolD_sound_of_HK}, \lean{Level.Unique}, \lean{protocolD_sound_unique}, \lean{protocolD_sound_concrete_unique_le}, \lean{protocolD_sound_concrete_ladder}, \lean{strategy_nonempty_concrete_ladder}, \lean{err_level0} & \lean{Soundness/Hypothesis.lean}\\
Theorem~\ref{thm:rbr-deep} for a schedule with later batches; the link from the model to the concrete protocol & by hand (\S\ref{sec:rbr}, \S\ref{sec:formal}) & \\
\addlinespace
Lemma~\ref{lem:near}, Theorem~\ref{thm:K} (descent step) & \lean{nearT_unique}, \lean{NearT.inList}, \lean{NearT.fold_eq}, \lean{NearT.degree_lt}, \lean{descent_step}, \lean{existsUnique_nearT} & \lean{Rounds/Extract.lean}\\
 & \lean{descend_iff_unique}, \lean{descend_step} (for one fold round and under (S2), the extraction predicate of the composed theorem is the statement about the unique near tuple) & \lean{Soundness/Extract.lean}\\
Theorem~\ref{thm:K} (a general predicate over a run; the decoder as an algorithm) & by hand & \\
Corollary~\ref{cor:K2} & \lean{k2_hfalse} (the predicate is admissible, given the pinned member); \lean{claimsHold_unique} (at most one such member); the rest by hand & \lean{Rounds/Batch.lean}, \lean{Soundness/Extract.lean}\\
Corollary~\ref{cor:johnson} & by hand (not connected to the composed theorem in Lean) & \\
\addlinespace
Theorem~\ref{thm:I} (tree bound) & \lean{tree_core}, \lean{interleave_listBound}, \lean{interleave_listBound_one}, \lean{interleave_listBound_same}, \lean{interleave_listBound_two_mul}, \lean{interleave_listBound_sq} & \lean{Interleave/Tree.lean}\\
 & \lean{interleave_listBound_closed}, \lean{interleave_listBound_choose} & \lean{Interleave/Closed.lean}\\
Proposition~\ref{prop:R2} (projection) & \lean{descent}, \lean{descent_div} & \lean{Interleave/Projection.lean}\\
 & \leansub{rs_descent}{1}, \leansub{rs_descent}{1}\lean{_div}, \lean{rs_descent} & \lean{Interleave/ReedSolomon.lean}\\
Proposition~\ref{prop:two} (two lanes) & \lean{listBound_prod}, \lean{listBound_prod_self}, \lean{listBound_prod_two_mul} & \lean{Interleave/TwoLanes.lean}\\
Hypothesis~\ref{hyp:HK} implies Hypothesis~\ref{hyp:LB} & \lean{rs_interleave}, \lean{rs_interleave_same}, \lean{rs_interleave_two_mul}, \lean{rs_listBound_min} & \lean{Interleave/ReedSolomon.lean}\\
 & \lean{card_fullT_filter_le}, \lean{lemmaF_exc_of_rsListBound}, \lean{lemmaF_exc_of_HK} & \lean{Interleave/Bridge.lean}\\
 & \lean{concrete_unique}, \lean{concrete_two_mul}, \lean{concrete_descent} (instances at a field of $2^{64}$ elements and a quadratic extension of it, with four lanes; not deployed parameters) & \lean{Interleave/Axioms.lean}\\
\addlinespace
Theorem~\ref{thm:quad}, Propositions~\ref{prop:affine}, \ref{prop:three}, \ref{prop:split} & by hand & \\
all bounds of \S\ref{sec:proved} for restricted classes of words & by hand, in the project's unpublished working notes; the proofs are not reproduced in this paper & \\
\addlinespace
Proposition~\ref{prop:params} (parameters) & \lean{level0_point_error}, \lean{point_error_level0_arith_11}, \lean{level0_ladder_error}, \lean{fold_error_arith}, \lean{fold_error_last_arith} & \lean{Protocol/Params.lean}\\
 & \lean{err_level0} (the same numbers as errors of the composed model) & \lean{Soundness/Hypothesis.lean}\\
Proposition~\ref{prop:floor} (floor) & by hand & \\
\bottomrule
\end{longtable}}

\subsection*{By file}

Every Lean file of the six libraries, with the statements of the paper it formalizes, except the six root files, which only import the others, and the files named \texttt{Axioms.lean}, which contain the axiom checks and instantiations at the deployed field sizes.

{\small
\begin{longtable}{@{}L{0.34\textwidth}L{0.62\textwidth}@{}}
\toprule
file under \texttt{lean/} & formalizes\\
\midrule
\endhead
\lean{DeepFold/Basic.lean} & Definition~\ref{def:cand} (the polynomial $\Delta$); Lemma~\ref{lem:L1}(c) and the degree of $\Delta$, one lane\\
\lean{DeepFold/Class.lean} & Lemma~\ref{lem:L1}(a), (b); Lemma~\ref{lem:L2}(a), (b), one lane\\
\lean{DeepFold/Agreement.lean} & Lemma~\ref{lem:L2}(c), one lane\\
\lean{DeepFold/Counting.lean} & Lemma~\ref{lem:L1}(d); Definition~\ref{def:exc}; Theorem~\ref{thm:count} in abstract form\\
\lean{DeepFold/Theorem.lean} & Theorem~\ref{thm:count} and the first half of Lemma~\ref{lem:F}(2) for a single lane\\
\lean{DeepFold/Lanes.lean} & Definition~\ref{def:ev}; Lemma~\ref{lem:E}; Lemmas~\ref{lem:L1}, \ref{lem:L2} and Theorem~\ref{thm:count} with lanes\\
\lean{DeepFold/Claim.lean} & Lemma~\ref{lem:F}(2) for any finite set of candidates\\
\lean{DeepFold/Full.lean} & Definition~\ref{def:cand} (the set $\cand$ is finite)\\
\lean{DeepFold/Main.lean} & Lemma~\ref{lem:F}\\
\addlinespace
\lean{Protocol/Potential.lean} & Lemma~\ref{lem:T} in abstract form (means, Markov, the potential of a process)\\
\lean{Protocol/Lists.lean} & Definition~\ref{def:lists} (lists) and its link with Definition~\ref{def:cand}\\
\lean{Protocol/Pinning.lean} & Lemma~\ref{lem:P}\\
\lean{Protocol/Level.lean} & Definition~\ref{def:ev} (the exponents $e_j$); Lemma~\ref{lem:E} (E2) and the bound behind (E4) for $e_j$; Lemma~\ref{lem:T} for one batch; row M0 of Theorem~\ref{thm:rbr-deep} for one batch; Hypothesis~\ref{hyp:LB} in its stated and diagonal forms\\
\lean{Protocol/Ladder.lean} & Lemma~\ref{lem:T} for a schedule; row M0 of Theorem~\ref{thm:rbr-deep} for a schedule\\
\lean{Protocol/Params.lean} & Proposition~\ref{prop:params}\\
\addlinespace
\lean{Rounds/Basic.lean} & Definition~\ref{def:fold}; Definition~\ref{def:pull}; the conditions $(\mathrm I_{i,j})$ and $(\mathrm I^B_{i,j})$ of \S\ref{sec:state} and \S\ref{sec:ks}; the updates \eqref{eq:claim-update} and \eqref{eq:carry}\\
\lean{Rounds/Batch.lean} & transition M1 of \S\ref{sec:transitions}; the predicate of Corollary~\ref{cor:K2}\\
\lean{Rounds/Fold.lean} & transition M2\\
\lean{Rounds/Query.lean} & transition M3, with the conversion \eqref{eq:convert}\\
\lean{Rounds/Final.lean} & transitions M4 and M5\\
\lean{Rounds/Extract.lean} & Lemma~\ref{lem:near}; the descent step of Theorem~\ref{thm:K}\\
\lean{Rounds/Pinned.lean} & transitions M1 and M3 with the separation hypothesis taken from Lemma~\ref{lem:P}\\
\addlinespace
\lean{Soundness/RBR.lean} & the definition of round-by-round soundness recalled in \S\ref{sec:annexb}; the sum of the errors bounds every strategy\\
\lean{Soundness/State.lean} & the state function of \S\ref{sec:state} and the transitions of \S\ref{sec:transitions} for the full state\\
\lean{Soundness/ProtocolD.lean} & Protocol~\ref{prot:D} (one batch per level, abstract model); Theorem~\ref{thm:rbr-deep} for it, in the form of \S\ref{sec:state}, ``commitment and opening as one protocol''\\
\lean{Soundness/Sound.lean} & Theorem~\ref{thm:rbr-deep} and Theorem~\ref{thm:K} as bounds on every prover strategy of the model; the existence of prover strategies (\S\ref{sec:composed})\\
\lean{Soundness/Extract.lean} & the descent of Theorem~\ref{thm:K} as unique decoding; uniqueness in Corollary~\ref{cor:K2}\\
\lean{Soundness/Hypothesis.lean} & Hypotheses~\ref{hyp:LB} and \ref{hyp:HK} as hypotheses of the composed theorem; the unique-decoding range with its two concrete instances; Proposition~\ref{prop:params} for the model\\
\addlinespace
\lean{Interleave/Basic.lean} & the setting of \S\ref{sec:interleave} (codes with an agreement bound, lists, interleaving)\\
\lean{Interleave/Tree.lean} & Theorem~\ref{thm:I}\\
\lean{Interleave/Closed.lean} & the closed forms of Theorem~\ref{thm:I}\\
\lean{Interleave/Projection.lean} & Proposition~\ref{prop:R2}\\
\lean{Interleave/Extension.lean} & Theorem~\ref{thm:I} for values in a finite-dimensional space (first paragraph of \S\ref{sec:interleave})\\
\lean{Interleave/ReedSolomon.lean} & Hypothesis~\ref{hyp:HK}; Theorem~\ref{thm:I} and Proposition~\ref{prop:R2} for Reed-Solomon codes; Hypothesis~\ref{hyp:HK} implies Hypothesis~\ref{hyp:LB}\\
\lean{Interleave/TwoLanes.lean} & Proposition~\ref{prop:two}\\
\lean{Interleave/Bridge.lean} & Lemma~\ref{lem:F}(1) from Hypothesis~\ref{hyp:HK}\\
\addlinespace
\lean{Constructions/Subspace.lean} & Proposition~\ref{prop:flats} (subspace polynomials are linearized)\\
\lean{Constructions/Flats.lean} & Proposition~\ref{prop:flats}\\
\lean{Constructions/Pole.lean} & Proposition~\ref{prop:pole}\\
\lean{Constructions/Padded.lean} & Proposition~\ref{prop:five8}(c)\\
\lean{Constructions/Gap.lean} & the degree count in the proof of Proposition~\ref{prop:five8}(c) (products of locators with a common top)\\
\lean{Constructions/UnionOfFlats.lean} & Proposition~\ref{prop:five8}(a) in its description by unions of flats (every member shares the top of the word)\\
\lean{Constructions/Glued.lean} & Proposition~\ref{prop:five8}(a), one list member per label set\\
\lean{Constructions/GluedCount.lean} & Proposition~\ref{prop:five8}(a), a lower bound on the count\\
\lean{Constructions/Concrete.lean} & Propositions~\ref{prop:flats}, \ref{prop:five8}, \ref{prop:pole} at the deployed sizes; Proposition~\ref{prop:pole} with the flats\\
\bottomrule
\end{longtable}}

\subsection*{The main statements, verbatim}

The following are copied from the Lean sources without change (Unicode symbols are typeset by the listing, which also breaks the lines that are longer than the page). A listing ends where the proof begins, and it does not repeat the \texttt{variable} declarations of its file, so some implicit arguments (the field and its instances, for instance) do not show.

% Generated from the Lean sources by copying each statement verbatim; do not edit by hand.
\medskip\noindent\textbf{\texttt{lemmaF\_exc}}, \texttt{lean/DeepFold/Main.lean}: Lemma~\ref{lem:F}(1).\par\nopagebreak
\begin{lstlisting}
theorem lemmaF_exc (D : Finset K) (O₀ O₁ : K → ι → K) (d A : ℕ) {ex : ι → ℕ}
    (hex : Function.Injective ex) {E : ℕ} (hE : ∀ u, ex u ≤ E) (Λ : ℕ)
    (hΛ : ∀ c : K × K, ((fullT D O₀ O₁ d A).filter (fun e => e.1 = c)).card ≤ Λ)
    (s : ℕ) {B : ℕ} (hB : 0 < B) :
    ((ExcL (fullT D O₀ O₁ d A) ex s B).card : ℚ) / (Fintype.card (Fin s → K × K) : ℚ) ≤
      (Λ : ℚ) ^ 3 * (Fintype.card K : ℚ) ^ 6 *
        (((d + 1 + E : ℕ) : ℚ) / Fintype.card K) ^ s / B
\end{lstlisting}

\medskip\noindent\textbf{\texttt{lemmaF\_round}}, \texttt{lean/DeepFold/Main.lean}: Lemma~\ref{lem:F}(2).\par\nopagebreak
\begin{lstlisting}
theorem lemmaF_round [CharP K 2] (D : Finset K) (O₀ O₁ : K → ι → K) (d A : ℕ) (ex : ι → ℕ)
    (w : ℕ) {s K₀ M : ℕ} (hMq : 2 * Fintype.card K ≤ M) (hMn : D.card.choose 2 ≤ M)
    {z : Fin s → K × K}
    (hz : z ∉ ExcL (fullT D O₀ O₁ d A) ex s (K₀ * (K₀ - 1) * (K₀ - M)))
    (W₀ W₁ : (ι → K[X]) →ₗ[K] K) (σ : K) (v : Fin s → K)
    (B : (ι → K[X]) → (ι → K[X]) → Prop)
    (hyp : ∀ F₀ G₀ : ι → K[X], (∀ u, (F₀ u).degree < (d : WithBot ℕ)) →
      (∀ u, (G₀ u).natDegree ≤ d) → A ≤ (agrPairL D O₀ O₁ F₀ G₀).card + 1 → B F₀ G₀ →
      (∀ i, (z i).1 ^ w * evL ex (z i) F₀ + evL ex (z i) G₀ = v i) → W₀ F₀ + W₁ G₀ ≠ σ)
    (vG vF σ₀ σ₁ : K) (φm : Fin s → K) :
    (univ.filter (fun c : K × K => ∃ P : ι → K[X],
      (∀ u, (P u).natDegree ≤ d) ∧ A ≤ (agrL D O₀ O₁ (c, P)).card ∧
      (∀ i, evL ex (z i) P = v i + ((z i).1 ^ w + c.1 + c.2 * (z i).2) * φm i) ∧
      ((1 + c.1) • W₀ + c.1 • W₁) P =
        vG + c.1 * (vG + σ) + c.1 ^ 2 * vF + c.2 * σ₀ + c.1 * c.2 * (σ₀ + σ₁) ∧
      (∀ F₀ G₀ : ι → K[X], (∀ u, (F₀ u).natDegree ≤ d) → (∀ u, (G₀ u).natDegree ≤ d) →
        (agrL D O₀ O₁ (c, P) \ agrPairL D O₀ O₁ F₀ G₀).card ≤ 1 → B F₀ G₀))).card < K₀
\end{lstlisting}

\medskip\noindent\textbf{\texttt{lemmaP}}, \texttt{lean/Protocol/Pinning.lean}: Lemma~\ref{lem:P}.\par\nopagebreak
\begin{lstlisting}
theorem lemmaP (D : Finset K) (O : K → ι → K) (d A : ℕ) {ex : ι → ℕ}
    (hex : Function.Injective ex) {E : ℕ} (hE : ∀ u, ex u ≤ E) {L : ℕ}
    (hL : (listW D O d A).card ≤ L) (s : ℕ) :
    dens (fun z : Fin s → K × K => z ∈ NonSep (listW D O d A) ex s) ≤
      (L.choose 2 : ℚ) * (((d + E : ℕ) : ℚ) / Fintype.card K) ^ s
\end{lstlisting}

\medskip\noindent\textbf{\texttt{levelPot\_step}}, \texttt{lean/Protocol/Level.lean}: Lemma~\ref{lem:T}(b).\par\nopagebreak
\begin{lstlisting}
theorem levelPot_step (hτ : 0 < τ) {z : Fin s → K × K} {k r : ℕ} {O : Word K (r + 1)}
    (h : levelPot τ D d thr s B z k (r + 1) O < 1) :
    z ∉ ExcL (fullT D (half₀ O) (half₁ O) d (thr (k + 1))) (exL (k + 1) r) s B ∧
      dens (fun c : K × K => 1 ≤ levelPot τ D d thr s B z (k + 1) r (foldW O c)) < τ
\end{lstlisting}

\medskip\noindent\textbf{\texttt{dens\_not\_good\_le}}, \texttt{lean/Protocol/Level.lean}: the point message of Theorem~\ref{thm:rbr-deep}, one batch.\par\nopagebreak
\begin{lstlisting}
theorem dens_not_good_le (hτ : 0 ≤ τ) (hB : 0 < B) {L ℓ : ℕ} (hℓ : 1 ≤ ℓ)
    (hL : ∀ k r, k + r = ℓ → ∀ O' : Word K r, (listW D O' d (thr k)).card ≤ L)
    (O : Word K ℓ) :
    dens (fun z : Fin s → K × K => ¬ Good τ D d thr s B ℓ O z) ≤
      ((L.choose 2 : ℚ) + (L : ℚ) ^ 3 * (Fintype.card K : ℚ) ^ 6 / B *
          ∑ j ∈ range ℓ, τ⁻¹ ^ j) * (((d + 2 ^ ℓ - 1 : ℕ) : ℚ) / Fintype.card K) ^ s
\end{lstlisting}

\medskip\noindent\textbf{\texttt{dens\_not\_goodS\_le}}, \texttt{lean/Protocol/Ladder.lean}: the point message of Theorem~\ref{thm:rbr-deep}, any schedule.\par\nopagebreak
\begin{lstlisting}
theorem dens_not_goodS_le (hτ : 0 ≤ τ) {L ℓ : ℕ} (hℓ : 1 ≤ ℓ)
    (hL : ∀ k r, k + r = ℓ → ∀ O' : Word K r, (listW D O' d (thr k)).card ≤ L)
    (O : Word K ℓ) (n : ℕ) :
    dens (fun p : Fin n → K × K => ¬ GoodS τ D d thr bt B ℓ O p) ≤
      (L.choose 2 : ℚ) * (((d + 2 ^ ℓ - 1 : ℕ) : ℚ) / Fintype.card K) ^ n +
        ∑ j ∈ range ℓ, τ⁻¹ ^ j * ((L : ℚ) ^ 3 * (Fintype.card K : ℚ) ^ 6 *
          (((d + 2 ^ ℓ - 1 : ℕ) : ℚ) / Fintype.card K) ^ held bt 0 n j / B)
\end{lstlisting}

\medskip\noindent\textbf{\texttt{level0\_point\_error}}, \texttt{lean/Protocol/Params.lean}: Proposition~\ref{prop:params}, twelve points at level 0.\par\nopagebreak
\begin{lstlisting}
theorem level0_point_error (hq : Fintype.card K = 2 ^ 128) (D : Finset K) (thr : ℕ → ℕ)
    (hL : ∀ k r, k + r = 6 → ∀ O' : Word K r,
      (listW D O' (2 ^ 24 + 5) (thr k)).card ≤ 2 ^ 40) (O : Word K 6) :
    dens (fun z : Fin 12 → K × K => ¬ Good (1 / 2 ^ 119) D (2 ^ 24 + 5) thr 12
      (2 ^ 137 * (2 ^ 137 - 1) * (2 ^ 137 - 2 ^ 129)) 6 O z) ≤ 1 / 2 ^ 110
\end{lstlisting}

\medskip\noindent\textbf{\texttt{Doomed}}, \texttt{lean/Rounds/Basic.lean}: the state $(\mathrm{I}^{B}_{i,j})$.\par\nopagebreak
\begin{lstlisting}
def Doomed (D : Finset K) (O : K → ι → K) (d A : ℕ) (ex : ι → ℕ) {s : ℕ}
    (z : Fin s → K × K) (v : Fin s → K) (W : (ι → K[X]) →ₗ[K] K) (σ : K)
    (B : (ι → K[X]) → Prop) : Prop :=
  ∀ U : ι → K[X], InList D O d A U → B U → (∀ m, evL ex (z m) U = v m) → W U ≠ σ
\end{lstlisting}

\medskip\noindent\textbf{\texttt{batch\_transition}}, \texttt{lean/Rounds/Batch.lean}: transition M1.\par\nopagebreak
\begin{lstlisting}
theorem batch_transition (D : Finset K) (O : K → ι → K) (d A : ℕ) (ex : ι → ℕ) {s : ℕ}
    (z : Fin s → K × K) (y : Fin s → K) {J : ℕ} (Wk : Fin J → (ι → K[X]) →ₗ[K] K)
    (ck : Fin J → K) (B : (ι → K[X]) → Prop)
    (hsep : ∀ U U' : ι → K[X], InList D O d A U → InList D O d A U' →
      (∀ m, evL ex (z m) U = y m) → (∀ m, evL ex (z m) U' = y m) → U = U')
    (hfalse : ∀ U : ι → K[X], InList D O d A U → B U → (∀ m, evL ex (z m) U = y m) →
      ∃ k, Wk k U ≠ ck k) :
    (univ.filter (fun l : K =>
      ¬ Doomed D O d A ex z y (batchW Wk l) (batchC ck l) B)).card ≤ J - 1
\end{lstlisting}

\medskip\noindent\textbf{\texttt{fold\_transition}}, \texttt{lean/Rounds/Fold.lean}: transition M2.\par\nopagebreak
\begin{lstlisting}
theorem fold_transition [CharP K 2] (D : Finset K) (O : K → ι ⊕ ι → K) (d A : ℕ) (ex : ι → ℕ)
    (w : ℕ) {s K₀ M : ℕ} (hMq : 2 * Fintype.card K ≤ M) (hMn : D.card.choose 2 ≤ M)
    {z : Fin s → K × K}
    (hz : z ∉ ExcL (fullT D (half₀ O) (half₁ O) d A) ex s (K₀ * (K₀ - 1) * (K₀ - M)))
    (W : (ι ⊕ ι → K[X]) →ₗ[K] K) (σ : K) (v : Fin s → K) (B : (ι ⊕ ι → K[X]) → Prop)
    (hI : Doomed D O d (A - 1) (exHalves w ex) z v W σ B)
    (vG vF σ₀ σ₁ : K) (φm : Fin s → K) :
    (univ.filter (fun c : K × K =>
      ¬ Doomed D (foldWord c O) d A ex z (foldVal w z v φm c) (foldWeight W c)
        (foldClaim σ vG vF σ₀ σ₁ c) (pull D O d A c B))).card < K₀
\end{lstlisting}

\medskip\noindent\textbf{\texttt{query\_transition}}, \texttt{lean/Rounds/Query.lean}: transition M3.\par\nopagebreak
\begin{lstlisting}
theorem query_transition [CharP K 2] (D : Finset K) (w : K → K) (d A : ℕ) {s : ℕ}
    (ζ v : Fin s → K) (W : K[X] →ₗ[K] K) (σ : K) (B : K[X] → Prop)
    (hI : Doomed₁ D w d A ζ v W σ B)
    (D' : Finset K) (O' : K → ι' → K) (d' A' : ℕ) (ex' : ι' → ℕ) {s' : ℕ}
    (z' : Fin s' → K × K) (y' : Fin s' → K)
    (hsep : ∀ U U' : ι' → K[X], InList D' O' d' A' U → InList D' O' d' A' U' →
      (∀ m, evL ex' (z' m) U = y' m) → (∀ m, evL ex' (z' m) U' = y' m) → U = U')
    (low : (ι' → K[X]) →ₗ[K] K[X]) (C : K[X])
    (hdeg : ∀ U : ι' → K[X], InList D' O' d' A' U → (low U + C).natDegree ≤ d)
    (hWC : W C = 0) (t : ℕ) :
    (((Fintype.piFinset (fun _ : Fin t => D)) ×ˢ (univ : Finset K)).filter
      (fun xl : (Fin t → K) × K =>
        ¬ Doomed D' O' d' A' ex' z' y' (nextW W low ζ xl.1 xl.2) (nextC σ C w ζ v xl.1 xl.2)
          (fun U => B (low U + C)))).card ≤
      (A - 1) ^ t * Fintype.card K + D.card ^ t * (s + t)
\end{lstlisting}

\medskip\noindent\textbf{\texttt{final\_transition}}, \texttt{lean/Rounds/Final.lean}: transition M4.\par\nopagebreak
\begin{lstlisting}
theorem final_transition (D : Finset K) (w : K → K) (d A : ℕ) {s : ℕ} (ζ v : Fin s → K)
    (W : K[X] →ₗ[K] K) (σ : K) (B : K[X] → Prop) (hI : Doomed₁ D w d A ζ v W σ B)
    (P : K[X]) (hP : P.natDegree ≤ d) (t : ℕ) :
    ((Fintype.piFinset (fun _ : Fin t => D)).filter (fun x =>
      (∀ m, P.eval (ζ m) = v m) ∧ (∀ q, P.eval (x q) = w (x q)) ∧ W P = σ ∧
        ¬ (A ≤ (D.filter (fun x => P.eval x = w x)).card ∧ ¬ B P))).card ≤ (A - 1) ^ t
\end{lstlisting}

\medskip\noindent\textbf{\texttt{descent\_step}}, \texttt{lean/Rounds/Extract.lean}: the descent step of Theorem~\ref{thm:K}.\par\nopagebreak
\begin{lstlisting}
theorem descent_step (D : Finset K) (O : K → ι ⊕ ι → K) (d A : ℕ) (c : K × K)
    (B : (ι ⊕ ι → K[X]) → Prop) (hdA : d + 3 ≤ A) {P : ι → K[X]}
    (hP : ∀ u, (P u).natDegree ≤ d) (hnb : ¬ pull D O d A c B P) :
    InList D (foldWord c O) d A P ∧
      ∃ T : ι ⊕ ι → K[X], NearT D O d (agr D (foldWord c O) P) T ∧ InList D O d (A - 1) T ∧
        ¬ B T ∧ P = foldTuple c T ∧
        ∀ T' : ι ⊕ ι → K[X], NearT D O d (agr D (foldWord c O) P) T' → T' = T
\end{lstlisting}

\medskip\noindent\textbf{\texttt{StateFunction}}, \texttt{lean/Soundness/RBR.lean}: the definition of a round-by-round state function, properties (i) to (iv).\par\nopagebreak
\begin{lstlisting}
structure StateFunction (Γ : PublicCoin.{u, v, w}) (start win : Γ.Pos → Prop) (err : Γ.Pos → ℚ)
    (doomed : Γ.Pos → Prop) (doomedP : ∀ s, Γ.Msg s → Prop) : Prop where
  /-- (i) -/
  start_doomed : ∀ s, start s → doomed s
  /-- (ii) -/
  prover_doomed : ∀ s (p : Γ.Msg s), doomed s → doomedP s p
  /-- (iii) -/
  verifier_doomed : ∀ s (p : Γ.Msg s), doomedP s p →
    dens (fun c : Γ.Chal s => ¬ doomed (Γ.next s p c)) ≤ err s
  /-- (iv) -/
  doomed_not_win : ∀ s, doomed s → ¬ win s
\end{lstlisting}

\medskip\noindent\textbf{\texttt{RBRSound}}, \texttt{lean/Soundness/RBR.lean}: round-by-round soundness.\par\nopagebreak
\begin{lstlisting}
def RBRSound (Γ : PublicCoin.{u, v, w}) (start win : Γ.Pos → Prop) (err : Γ.Pos → ℚ) : Prop :=
  ∃ (doomed : Γ.Pos → Prop) (doomedP : ∀ s, Γ.Msg s → Prop),
    StateFunction Γ start win err doomed doomedP
\end{lstlisting}

\medskip\noindent\textbf{\texttt{Level.OK}}, \texttt{lean/Soundness/ProtocolD.lean}: the standing assumptions and Hypothesis~\ref{hyp:LB} (diagonal form) for a level of the model.\par\nopagebreak
\begin{lstlisting}
structure Level.OK [Fintype K] (M : ℕ) (lv : Level K) : Prop where
  /-- `ℓ ≥ 1`. -/
  ell_pos : 1 ≤ lv.ℓ
  /-- The domain is not empty. -/
  nonempty : lv.D.Nonempty
  /-- (S3). -/
  choose_le : lv.D.card.choose 2 ≤ M
  /-- `A^{(i,k)} + 1 ≤ A^{(i,k+1)}`. -/
  thr_step : ∀ k, lv.thr k + 1 ≤ lv.thr (k + 1)
  /-- `ListBound`. -/
  listBound : ∀ k r, k + r = lv.ℓ → ∀ O' : Word K r, (listW lv.D O' lv.d (lv.thr k)).card ≤ lv.L
\end{lstlisting}

\medskip\noindent\textbf{\texttt{win}}, \texttt{lean/Soundness/ProtocolD.lean}: the winning condition.\par\nopagebreak
\begin{lstlisting}
def win : Pos K → Prop
  | .done acc ext tru => acc ∧ ¬ (ext ∧ tru)
  | _ => False
\end{lstlisting}

\medskip\noindent\textbf{\texttt{protocolD\_rbr}}, \texttt{lean/Soundness/ProtocolD.lean}: Theorem~\ref{thm:rbr-deep} for the abstract model, one batch per level.\par\nopagebreak
\begin{lstlisting}
theorem protocolD_rbr [CharP K 2] (hτ : 0 < τ) (hMq : 2 * Fintype.card K ≤ M)
    (hΘ : 0 < tripleThr K₀ M) (J : ℕ) :
    RBRSound (protocolD K T J) (start M) win (err τ K₀ M J)
\end{lstlisting}

\medskip\noindent\textbf{\texttt{protocolD\_sound}}, \texttt{lean/Soundness/Sound.lean}: every prover strategy wins with probability at most the sum of the errors.\par\nopagebreak
\begin{lstlisting}
theorem protocolD_sound [CharP K 2] (hτ : 0 < τ) (hMq : 2 * Fintype.card K ≤ M)
    (hΘ : 0 < tripleThr K₀ M) (J : ℕ) (lv : Level K) (rest : List (Level K))
    (hOK : ∀ l ∈ lv :: rest, l.OK M) (n : ℕ)
    (σ : (protocolD K T J).Strategy n (Pos.commit lv rest)) :
    (protocolD K T J).winProb win n (Pos.commit lv rest) σ ≤
      budget T τ K₀ M J (Pos.commit lv rest)
\end{lstlisting}

\medskip\noindent\textbf{\texttt{protocolD\_sound\_extract}}, \texttt{lean/Soundness/Sound.lean}: Theorem~\ref{thm:K} composed, canonical predicate.\par\nopagebreak
\begin{lstlisting}
theorem protocolD_sound_extract [CharP K 2] (hτ : 0 < τ) (hMq : 2 * Fintype.card K ≤ M)
    (hΘ : 0 < tripleThr K₀ M) (J : ℕ) (lv : Level K) (rest : List (Level K))
    (hOK : ∀ l ∈ lv :: rest, l.OK M) (n : ℕ)
    (σ : (protocolD K T J).Strategy n (Pos.commit lv rest)) :
    (protocolD K T J).winProb winExtract n (Pos.commit lv rest) σ ≤
      budget T τ K₀ M J (Pos.commit lv rest)
\end{lstlisting}

\medskip\noindent\textbf{\texttt{Terminal.AdmitsZero}}, \texttt{lean/Soundness/Sound.lean}: the condition on the tail under which the prover has a final message.\par\nopagebreak
\begin{lstlisting}
def Terminal.AdmitsZero : Prop :=
  ∀ W : (Lane 0 → K[X]) →ₗ[K] K, ∑ x, T.wtab W 0 x * T.ftab W 0 x = 0
\end{lstlisting}

\medskip\noindent\textbf{\texttt{Level.Fits}}, \texttt{lean/Soundness/Sound.lean}: the condition on two consecutive levels under which the prover has a high part.\par\nopagebreak
\begin{lstlisting}
def Level.Fits (lv lv' : Level K) : Prop :=
  ∀ U : Lane lv'.ℓ → K[X], (∀ u, (U u).natDegree ≤ lv'.d) → (lv'.low U).natDegree ≤ lv.d
\end{lstlisting}

\medskip\noindent\textbf{\texttt{ladderFits}}, \texttt{lean/Soundness/Sound.lean}: the same for a ladder.\par\nopagebreak
\begin{lstlisting}
def ladderFits : Level K → List (Level K) → Prop
  | _, [] => True
  | lv, lv' :: rest => lv.Fits lv' ∧ ladderFits lv' rest
\end{lstlisting}

\medskip\noindent\textbf{\texttt{strategy\_nonempty\_ladder}}, \texttt{lean/Soundness/Sound.lean}: prover strategies exist on every such ladder with every such tail.\par\nopagebreak
\begin{lstlisting}
theorem strategy_nonempty_ladder (hT : T.AdmitsZero) (J : ℕ) (lv : Level K)
    (rest : List (Level K)) (h : ladderFits lv rest) (n : ℕ) :
    Nonempty ((protocolD K T J).Strategy n (Pos.commit lv rest))
\end{lstlisting}

\medskip\noindent\textbf{\texttt{protocolD\_sound\_of\_HK}}, \texttt{lean/Soundness/Hypothesis.lean}: the same from Hypothesis~\ref{hyp:HK}.\par\nopagebreak
\begin{lstlisting}
theorem protocolD_sound_of_HK [Fintype F] [CharP K 2] (hτ : 0 < τ) (hMq : 2 * Fintype.card K ≤ M)
    (hΘ : 0 < tripleThr K₀ M) (J : ℕ) (lv : Level K) (rest : List (Level K))
    (hHK : ∀ l ∈ lv :: rest, l.HK F M) (n : ℕ)
    (σ : (protocolD K T J).Strategy n (Pos.commit lv rest)) :
    (protocolD K T J).winProb win n (Pos.commit lv rest) σ ≤
      budget T τ K₀ M J (Pos.commit lv rest)
\end{lstlisting}

\medskip\noindent\textbf{\texttt{protocolD\_sound\_concrete\_unique\_le}}, \texttt{lean/Soundness/Hypothesis.lean}: a concrete instance with no hypothesis (one level, no tail; not a deployed parameter set).\par\nopagebreak
\begin{lstlisting}
theorem protocolD_sound_concrete_unique_le (n : ℕ)
    (σ : (protocolD E128 (Terminal.direct E128) 512).Strategy n (Pos.commit levelConcrete [])) :
    (protocolD E128 (Terminal.direct E128) 512).winProb win n (Pos.commit levelConcrete []) σ ≤
      1 / 2 ^ 109
\end{lstlisting}

\medskip\noindent\textbf{\texttt{protocolD\_sound\_concrete\_ladder}}, \texttt{lean/Soundness/Hypothesis.lean}: a concrete instance with no hypothesis on two levels with a tail of two rounds (not a deployed parameter set).\par\nopagebreak
\begin{lstlisting}
theorem protocolD_sound_concrete_ladder (n : ℕ)
    (σ : (protocolD E128 (Terminal.point E128 2) 512).Strategy n
      (Pos.commit levelConcrete [levelConcrete₁])) :
    (protocolD E128 (Terminal.point E128 2) 512).winProb win n
        (Pos.commit levelConcrete [levelConcrete₁]) σ ≤
      budget (Terminal.point E128 2) (1 / 2 ^ 119) (2 ^ 137) (2 ^ 129) 512
        (Pos.commit levelConcrete [levelConcrete₁])
\end{lstlisting}

\medskip\noindent\textbf{\texttt{interleave\_listBound}}, \texttt{lean/Interleave/Tree.lean}: Theorem~\ref{thm:I}.\par\nopagebreak
\begin{lstlisting}
theorem interleave_listBound (hC : AgreeLE D C d) (hda : d < a) (hL : ListBound D C a L)
    (m : Type*) [Finite m] :
    ListBound D (interleave m C) a (treeN L (a - d) (D.card - d))
\end{lstlisting}

\medskip\noindent\textbf{\texttt{rs\_descent}\textsubscript{\texttt{1}}}, \texttt{lean/Interleave/ReedSolomon.lean}: Proposition~\ref{prop:R2} from $\HK$.\par\nopagebreak
\begin{lstlisting}
theorem rs_descent₁ (hL : RSListBound₁ D d a L) (h2 : 2 * L ≤ Fintype.card K)
    (w : K → lanes → F) (S : Finset (lanes → F[X]))
    (hS : ∀ P ∈ S, (∀ u, (P u).natDegree ≤ d) ∧
      a ≤ (agr D w (fun x u => (P u).eval (algebraMap K F x))).card) :
    S.card ≤ 2 * L - 1 ∧ S.card * (Fintype.card K - L) ≤ L * Fintype.card K
\end{lstlisting}

\medskip\noindent\textbf{\texttt{lemmaF\_exc\_of\_HK}}, \texttt{lean/Interleave/Bridge.lean}: Lemma~\ref{lem:F}(1) from $\HK$.\par\nopagebreak
\begin{lstlisting}
theorem lemmaF_exc_of_HK (D₀ : Finset K) (O₀ O₁ : E → ι → E) (d A : ℕ) {ex : ι → ℕ}
    (hex : Function.Injective ex) {Eb : ℕ} (hE : ∀ u, ex u ≤ Eb) (L : ℕ)
    (hL : RSListBound₁ D₀ d A L) (hda : d < A) (h4 : D₀.card + 3 * d < 4 * A) (s : ℕ) {B : ℕ}
    (hB : 0 < B) :
    ((DeepFold.ExcL (DeepFold.fullT (D₀.image (algebraMap K E)) O₀ O₁ d A) ex s B).card : ℚ) /
        (Fintype.card (Fin s → E × E) : ℚ) ≤
      ((2 * L - 1 : ℕ) : ℚ) ^ 3 * (Fintype.card E : ℚ) ^ 6 *
        (((d + 1 + Eb : ℕ) : ℚ) / Fintype.card E) ^ s / B
\end{lstlisting}

\medskip\noindent\textbf{\texttt{flats\_quarter}}, \texttt{lean/Constructions/Flats.lean}: Proposition~\ref{prop:flats}, rate $1/4$.\par\nopagebreak
\begin{lstlisting}
theorem flats_quarter (D : AddSubgroup K) [Finite D] {s : ℕ} (hs : 1 ≤ s)
    (hD : Nat.card D = 2 ^ (s + 1)) :
    ((Nat.card D - 1 : ℕ) : ℕ∞) ≤
      (RSList (pts D) (fun x => x ^ (2 ^ s - 1)) (2 ^ (s - 1)) (2 ^ s - 1)).encard
\end{lstlisting}

\medskip\noindent\textbf{\texttt{list\_gives\_bad\_challenges}}, \texttt{lean/Constructions/Pole.lean}: Proposition~\ref{prop:pole}.\par\nopagebreak
\begin{lstlisting}
theorem list_gives_bad_challenges [Fintype K] (D : Finset K) (w : K → K) {d T : ℕ}
    (𝒞 : Finset K[X]) (h𝒞 : ∀ Ci ∈ 𝒞, Ci.natDegree ≤ d + 1 ∧ T ≤ (agree D w Ci).card)
    (hK : (d + 1) * 𝒞.card.choose 2 + D.card < Fintype.card K) :
    ∃ β₀ : K, β₀ ∉ D ∧
      𝒞.card * Fintype.card K ≤
        (univ.filter (fun c : K × K =>
          Close D (affFold (poleWord β₀) (poleNum w β₀) c) d T)).card
\end{lstlisting}

\medskip\noindent\textbf{\texttt{five\_eighths\_slack}}, \texttt{lean/Constructions/Padded.lean}: Proposition~\ref{prop:five8}(c).\par\nopagebreak
\begin{lstlisting}
theorem five_eighths_slack (D H : AddSubgroup K) [Finite D] (hHD : H ≤ D)
    (hH : 2 * Nat.card H = Nat.card D) {m : ℕ} (hD : Nat.card D = 8 * m) (P : Finset K)
    (hP : ∀ x ∈ P, x ∈ H) (hPc : P.card = 3 * m) :
    ((Nat.card D - 2 : ℕ) : ℕ∞) ≤
      (RSList (pts D) (padWord P (2 * m)) (4 * m + 1) (5 * m)).encard
\end{lstlisting}

\medskip\noindent\textbf{\texttt{glued\_list\_card}}, \texttt{lean/Constructions/GluedCount.lean}: Proposition~\ref{prop:five8}(a), a lower bound on its count.\par\nopagebreak
\begin{lstlisting}
theorem glued_list_card (D U : AddSubgroup K) [Finite D] (hUD : U ≤ D)
    (hD : Nat.card D = 8 * Nat.card U) :
    ((Nat.card U - 1 : ℕ) : ℕ∞) ≤
      (RSList (pts D) (fun x => (subPoly U).eval x ^ 5)
        (4 * Nat.card U + 1) (5 * Nat.card U)).encard
\end{lstlisting}

